\documentclass[tikz,border=0pt]{standalone}
\usepackage{tikz}
\usetikzlibrary{shadings,fadings}

% 光晕：中心不透明，边缘透明
\tikzfading[name=glow, inner color=transparent!0, outer color=transparent!100]
% 暗角：中心透明，边缘不透明
\tikzfading[name=vignette, inner color=transparent!100, outer color=transparent!0]

\begin{document}
	\begin{tikzpicture}[line cap=round]
		\def\W{16}      % 画布宽
		\def\H{9}       % 画布高
		\def\HZ{3}      % 地平线高度

		\clip (0,0) rectangle (\W,\H);

		% ===== 1. 天空渐变 =====
		\shade[top color=blue!25!black, bottom color=magenta!40!black]
		(0,\HZ) rectangle (\W,\H);

		% ===== 2. 地面渐变 =====
		\shade[top color=magenta!20!black, bottom color=black]
		(0,0) rectangle (\W,\HZ);

		% ===== 3. 星空 =====
		\pgfmathsetseed{42}
		\foreach \i in {1,...,160} {%
			\pgfmathsetmacro{\sx}{\W*rnd}%
			\pgfmathsetmacro{\sy}{\HZ + 5.5*rnd}%
			\pgfmathsetmacro{\sr}{0.2 + 0.7*rnd}%
			\pgfmathsetmacro{\so}{0.15 + 0.75*rnd}%
			\fill[white, opacity=\so] (\sx,\sy) circle (\sr pt);
		}

		% ===== 4. 太阳外发光 =====
		\fill[orange!60!red, path fading=glow, opacity=0.6]
		(8,4.8) circle (2.6);

		% ===== 5. 太阳本体（带条纹）=====
		\begin{scope}
			\clip (8,4.8) circle (1.5);
			\shade[top color=yellow!95!white, bottom color=red!65!magenta]
			(6.5,3.3) rectangle (9.5,6.3);
			% 水平条纹，越往下越粗
			\foreach \yy/\hh in {%
				4.25/0.06, 4.13/0.07, 4.00/0.08, 3.86/0.09,
				3.71/0.10, 3.55/0.11, 3.38/0.12%
			} {%
				\fill[black!85!blue] (6.5,\yy) rectangle (9.5,\yy-\hh);
			}
		\end{scope}

		% ===== 6. 地平线光晕 =====
		\fill[magenta!55!white, path fading=glow, opacity=0.5]
		(8,\HZ) ellipse [x radius=8.5, y radius=0.7];

		% ===== 7. 远山剪影 =====
		\fill[black!90!blue]
		(0,3) --
		(0,3.2) -- (0.8,3.55) -- (1.6,3.25) -- (2.5,3.75) --
		(3.4,3.3) -- (4.2,3.6) -- (5.1,3.15) -- (6.0,3.5) --
		(7.0,3.25) -- (7.8,3.7) -- (8.6,3.3) -- (9.4,3.6) --
		(10.3,3.2) -- (11.2,3.65) -- (12.1,3.3) -- (13.0,3.55) --
		(14.0,3.2) -- (15.0,3.5) -- (16,3.25) --
		(16,3) -- cycle;

		% ===== 8. 透视网格 =====
		% 垂直线：全部汇聚到消失点 (8,3)
		\foreach \x in {-30,-28,...,46} {%
			\draw[cyan!80!white, opacity=0.28, line width=0.4pt]
			(\x,0) -- (8,\HZ);
		}
		% 水平线：越靠近地平线越密
		\foreach \k in {1,...,32} {%
			\pgfmathsetmacro{\gy}{\HZ - \HZ*(\k/32)*(\k/32)}%
			\draw[cyan!80!white, opacity=0.20, line width=0.35pt]
			(0,\gy) -- (\W,\gy);
		}

		% ===== 9. 边缘暗角 =====
		\fill[black, path fading=vignette, opacity=0.75]
		(0,0) rectangle (\W,\H);

	\end{tikzpicture}
\end{document}